% Generated from docs/REPORT_DRAFT.md; edit the manuscript then rerun the converter.
\chapter{Giới thiệu}

\section{Bối cảnh}
\reportfigure{movie_kg_overview.pdf}{Khái quát miền dữ liệu Movie Knowledge Graph}{fig:movie-kg-overview}

Một bộ phim không tồn tại như một bản ghi độc lập. Nó liên hệ nhiều–nhiều với
diễn viên, đạo diễn, thể loại, từ khóa và hãng sản xuất. Một người có thể đồng
thời là diễn viên và đạo diễn; hai diễn viên có thể đóng chung trong nhiều phim;
một nhóm phim có thể liên hệ gián tiếp qua đạo diễn, cast hoặc chủ đề. Dữ liệu
này lại nằm ở nhiều nguồn với định danh, mức đầy đủ và cách chấm điểm khác nhau.
TMDB cung cấp metadata, credits, keywords và external ID, trong khi IMDb cung cấp
rating và vote count độc lập.

Các câu hỏi như “đường liên hệ giữa hai diễn viên là gì?”, “đạo diễn nào thường
làm phim thuộc một thể loại?” hoặc “phim nào tương tự và vì sao?” không chỉ là
lookup một bản ghi. Chúng yêu cầu duyệt nhiều quan hệ và giữ lại bằng chứng của
đường duyệt. Knowledge Graph phù hợp với bài toán này vì thực thể và quan hệ được
biểu diễn tường minh, còn câu hỏi nghiệp vụ có thể ánh xạ thành graph pattern.

\section{Phát biểu bài toán}

Đề tài xây dựng một biểu diễn thống nhất cho năm loại thực thể \texttt{Movie}, \texttt{Person},
\texttt{Genre}, \texttt{Keyword} và \texttt{Studio}; tích hợp chính xác dữ liệu TMDB–IMDb; cung cấp
truy vấn nhiều bước và suy diễn có thể kiểm chứng; đồng thời minh họa giá trị của
graph thông qua hỏi–đáp và gợi ý phim có giải thích.

Bài toán không chỉ là “đưa dữ liệu vào Neo4j”. Một hệ thống đạt yêu cầu phải giải
quyết đồng thời:

\begin{enumerate}
\item Định danh: không dùng tên làm khóa vì tên có thể trùng hoặc thay đổi.
\item Provenance: phân biệt rating TMDB và IMDb, truy vết fact gốc hoặc fact suy ra.
\item Chất lượng: phát hiện bản ghi lỗi, duplicate, orphan và quan hệ sai kiểu/hướng.
\item An toàn truy vấn: không ghép chuỗi người dùng vào Cypher.
\item Tái lập: cache, checksum, manifest, cấu hình và kết quả đo phải kiểm chứng được.
\item Trung thực thực nghiệm: phân biệt smoke test, silver evaluation và human review.
\end{enumerate}

\section{Câu hỏi nghiên cứu}

\textbf{RQ chính:} Knowledge Graph có thể tích hợp dữ liệu phim đa nguồn, hỗ trợ truy
vấn multi-hop và suy diễn, đồng thời tạo câu trả lời và gợi ý có bằng chứng như
thế nào?

Các câu hỏi phụ gồm:

\begin{itemize}
\item \textbf{RQ1:} Exact-ID enrichment và stable source ID bảo đảm identity/provenance đến mức nào?
\item \textbf{RQ2:} Property Graph hỗ trợ các competency question đa bước ra sao?
\item \textbf{RQ3:} Rule materialization và semantic entailment sinh tri thức mới như thế nào, và làm sao kiểm tra tri thức đó?
\item \textbf{RQ4:} IDF-weighted graph similarity cung cấp recommendation và explanation như thế nào so với các phương pháp overlap khác?
\item \textbf{RQ5:} Hệ thống có giới hạn gì về coverage, accuracy, validity và scalability?
\end{itemize}

\section{Mục tiêu}

Mục tiêu tổng quát là xây dựng một Movie Knowledge Graph chạy được từ dữ liệu
nguồn đến ứng dụng. Các mục tiêu kiểm chứng được là:

\begin{itemize}
\item Thu thập và chuẩn hóa 2.000–5.000 phim cùng thực thể liên quan.
\item Bảo toàn source ID, provenance và metadata quan trọng trên node/edge.
\item Nạp graph idempotent vào Neo4j, có constraint/index và validation.
\item Xây ontology RDF/OWL, RDF export và ít nhất 10 SPARQL query.
\item Xây ít nhất 10 Cypher query, gồm multi-hop, aggregation và shortest path.
\item Materialize \texttt{CO\_STARRED\_WITH} với evidence và chạy semantic entailment.
\item Cung cấp API/UI cho QA và explainable recommendation.
\item Đánh giá data quality, entity resolution, reasoning, QA, recommendation và latency.
\end{itemize}

\section{Phạm vi}

TMDB là nguồn graph chính; IMDb chỉ enrichment cho Movie bằng exact \texttt{imdb\_id}.
Mỗi phim giữ tối đa 20 cast member nhằm giới hạn kích thước và thời gian thu thập.
MVP không bao gồm Award/Wikidata, NLP trên overview, user-history personalization,
vector search, embedding, GraphRAG hoặc LLM-to-Cypher tự do. Các thành phần này là
hướng mở rộng, không được dùng để mô tả chức năng hiện có.

\section{Đóng góp}

Đóng góp của đề tài gồm:

\begin{itemize}
\item Pipeline đa nguồn storage-bounded và tái lập được.
\item Chiến lược identity dựa trên stable source ID thay vì tên.
\item Mô hình kết hợp Neo4j operational graph và RDF/OWL standards view.
\item Rule suy diễn lưu evidence, cùng semantic materializer/validator chạy được.
\item QA có LLM planner nhưng execution surface được kiểm soát.
\item Recommendation graph-native với explanation từ chính feature đóng góp.
\item Bộ test/evaluation phân loại rõ evidence và giới hạn validity.
\end{itemize}

\section{Cấu trúc báo cáo}

Chương 2 trình bày cơ sở lý thuyết; Chương 3 tổng hợp công trình liên quan; Chương
4 phân tích yêu cầu; Chương 5 mô tả ontology và graph schema; Chương 6 trình bày
kiến trúc và pipeline; Chương 7 mô tả truy vấn/suy diễn; Chương 8 trình bày hai
ứng dụng; Chương 9 báo cáo thực nghiệm; Chương 10 kết luận và hướng phát triển.
